\chapter{第十七讲\quad 分论概说}
\label{lecture:17}
\noindent{\small\hyperref[ref:lecture:17]{查阅本讲扩展规则}}\par

\section{第一节\quad 罪名与罪状}
\paragraph{构成判断和罪名输出分两步}
令 \(B\) 为共用要件，\(T_i\) 为不同的可选行为或对象。先用 \(E=B\land\bigvee_i T_i\) 检验构成；再用司法解释确定的命名规则 \(N(E,\{i:T_i\})\) 输出罪名。选择性罪名允许按已证实分支输出或合并表述，多个分支满足不自动生成数罪；行为单位与罪数仍依第十三讲另审。概括罪名用同一个正式罪名覆盖各法定方式，例如信用卡诈骗罪（第196条），不能为“恶意透支”造一个独立罪名。单一罪名如第232条故意杀人罪，也须保持统一名称。具体名称以相应司法解释及其补充规定为准。

罪状描述的四种结构，区别在于要件信息从何处补全：
\begin{enumerate}
 \item 简单罪状以较少文字标识行为，如第232条；结合体系和合法解释明确其要件。
 \item 叙明罪状在条文中展开主要行为、主体、结果或目的。
 \item 引证罪状读取刑法内部所指条款，再代入本款增减的要件。如第115条第二款引用前款所涉行为、结果，再适用本款的过失要求。
 \item 空白罪状依法读取外部规范补充特定谓词。如第345条第二款的滥伐林木相关判断须查森林法规则。外部违规只补足对应谓词，数额、数量等入罪条件另行证明。
\end{enumerate}

\paragraph{例题推导}
甲以出卖为目的拐走十五岁女子，误以为是男童。设题中二种对象皆在第240条共同保护范围，主观上对拐卖行为有认识及出卖目的，客观上为拐卖妇女。先检验 \(B\) 成立，再验证客观分支 \(T_{\text{妇女}}\)；在教材关于选择性对象错误的解释下，选择对象之间的误认不阻却同一法定构成内的故意，按已完成的客观分支定拐卖妇女罪既遂。若甲以为对象是成年男子，且没有证据证明其认识到了法定保护对象的可能性，则不能直接套用该同一法定构成内的错误模型，应另查认识与构成。罪名表述能随分支变化，不能反过来替代主客观要件审查。

\section{第二节\quad 注意规定与法律拟制}
\paragraph{删除测试及其证明负担}
固定待考察的事实域 \(D\)，令 \(\mathcal R_0\) 为保留一般规定后的规则集，\(r\) 为待检验条款，\(Q(f)\) 为目标罪名结论。教材的删除测试分两类：
\[
\begin{aligned}
&\forall f\in D,\quad
  f\vdash_{\mathcal R_0\cup\{r\}}Q(f)
       \ \Longrightarrow\ f\vdash_{\mathcal R_0}Q(f)
       &&\text{结论可以由一般规则复现};\\
&\exists f\in D,\quad
  f\vdash_{\mathcal R_0\cup\{r\}}Q(f),\quad
  f\nvdash_{\mathcal R_0}Q(f)
       &&\text{该分支的结论依赖特别条款}.
\end{aligned}
\]
第一类在讨论范围内符合注意规定的结构。第二类中，若特别条款明确将原本不满足该罪构成的情形按该罪处罚，则符合教材所称法律拟制结构。判断第二类应指出一般构成的具体缺项，并由题设证明该缺项，不能把暂未查清的事实当成缺项。对同一条文的不同事实分支分别检测；拟制只在明文适用域内运作，不能类推。

\paragraph{例题推导}
甲盗窃完成后当场为抗拒抓捕而以暴力相威胁。题设前置盗窃、当场性、抗拒抓捕的目的以及暴力威胁均已证明；该威胁发生于取财之后，只为抗拒抓捕，缺乏以暴力夺取财物的普通抢劫结构。令 \(r\) 为第269条，保留 \(r\) 时，各要件合取推出依第263条按抢劫罪论处；移除 \(r\) 后，普通抢劫结构缺项，仅有这些事实无法推出普通抢劫，故转化结论依赖拟制条款。反例是甲次日才在异地威胁追来的人，“当场”失败，不能依据第269条转化，后续行为应独立检验。

\paragraph{一张两列对照卡}
第287条所举的利用计算机实施盗窃，如果已经具备盗窃全部要件，去掉该条仍由一般盗窃条款推出结论，属于教材的注意规定例。第267条第二款的携带凶器抢夺，若只有携带而未向被害人施加压制性暴力、胁迫，普通抢劫要件并未满足，但明文赋予抢劫的处理效果。教材把第238、247、248、292条部分致伤致死情形放在拟制框架下理解，尤其涉及过失致结果、行为类型和共犯归责时，要明示所采用的解释；不能把第269条的拟制规则类推到这些条款或倒推每位参与者均负转化后的责任。第196条第三款的盗窃信用卡并使用同样应先检验明文要件，再根据使用场景讨论删除测试，不宜凭罪名标签省略行为链。

\section*{记忆与闭卷复述}
记“罪状填空，罪名映射；删条重推，查缺哪项”。复述时先写共同要件和可选要件，随后给出法定名称。对注意规定与拟制，用同一组事实分别在 \(\mathcal R_0\) 和 \(\mathcal R_0\cup\{r\}\) 下推导，逐项圈出哪一个一般要件缺失，再换一项关键事实检验边界。
\begin{enumerate}
 \item \textbf{自测。}某条款要求违反特定外部规范，现已证明行为违反该规范，但法定数量门槛未达。能否由违规推出犯罪？\quad\textbf{核对。}不能。空白罪状引用只补足一项，其他并列构成要件仍须成立。
 \item \textbf{自测。}甲用计算机实施已满足全部要件的盗窃；乙携带凶器抢夺，未展示凶器或压制被害人。删除第287条、第267条第二款之后各有什么差异？\quad\textbf{核对。}甲的盗窃结论可依一般盗窃规则复现；乙在给定事实下的抢劫处理依赖第267条第二款，仍须单独评价其抢夺事实。
\end{enumerate}
